\chapter{Cognitive Task Analysis and Algebraic Transfer}
\label{sec:p19-task-analysis-algebra}

A learner who solves one problem but not a mathematically equivalent
one may be missing a representational skill, selecting a different
procedure, or misunderstanding the requested task. An instructional
designer therefore needs more than a topic label such as algebra,
problem solving, or critical thinking. The question is which knowledge
and operations the learner can use, under which conditions, and what
a proposed intervention changes.

Koedinger and Roll's 2011 typeset presentation of \emph{Learning to Think:
Cognitive Mechanisms of Knowledge Transfer} develops a component-based
account of this problem \cite[galley pp.~699--716]{KoedingerRoll2011TransferTypeset}.
That presentation, with its own pagination, is kept distinct from the
published Oxford chapter. Three original research papers supply the
empirical records below: Koedinger--Nathan's 2004 difficulty-factor
assessments, Koedinger--Alibali--Nathan's 2008 representation-complexity
comparisons, and Koedinger--McLaughlin's 2010 substitution intervention
\cite{KoedingerNathan2004StoryCTA,KoedingerAlibaliNathan2008TradeoffsCTA,
KoedingerMcLaughlin2010LanguageTransfer}.
Their different participants, tasks, outcome definitions and designs are
not combined into a single demonstration of general learning to think.
\index{cognitive task analysis}\index{knowledge components}
\index{algebra!representational transfer}\index{difficulty factors assessment}

\section{Knowledge components and the scope of transfer}
\label{sec:p19-cta-components}

\begin{hbdefinition}{def:p19-cta-components}{A component account of transfer}
In the selected account, a \emph{knowledge component} is a proposed
reusable piece of knowledge or mental function that participates in
performing a task. Components can include procedures, schemas,
representations, retrieval conditions, metacognitive strategies,
learning dispositions and other relevant acquired changes. A component
is not identified merely by a question's topic, a word in a curriculum,
or a correct answer.

Its \emph{scope} is the set of situations in which it can appropriately
be used. Broad verbal wording does not establish broad applicability.
A proposed component model must specify the common operation or
representation, its application conditions, and the task-specific
requirements that are not shared.
This is an empirical cognitive-modeling account, not a deductive axiom
that every observed transfer has a uniquely identified component
\cite[galley pp.~700--704,710--713]{KoedingerRoll2011TransferTypeset}.
\end{hbdefinition}

The source contrasts the size or sophistication of an idea with the
extent of its productive reuse. A comparatively local capability can
support many later tasks: decoding a new written word uses acquired
letter--sound relations, while distribution in algebra supports
transformations across many mathematical contexts. Decoding alone does
not establish comprehension, and being able to recite a distribution
formula does not establish its correct application in an unfamiliar
expression. The term \emph{small-big idea} concerns usefulness across
contexts, not a new class of formal logical principles.

The account distinguishes three limitations of indiscriminately general
instruction. A heuristic may omit the conditions under which to use it;
a learner may know a strategy but lack the domain knowledge needed to
carry it out; or the proposed general capability may already be available,
so teaching its name does not repair the actual bottleneck. These are
different explanations for a failed instructional contrast
\cite[galley pp.~701--703]{KoedingerRoll2011TransferTypeset}.

\begin{hbexample}{def:p19-cta-like-terms}{Application conditions in combining terms}
Use ordinary real arithmetic, with \(x\) the real input. The expressions
\(3x+4x\), \(x+5x\), and \(-3x+4x\) are respectively
\(7x\), \(6x\), and \(x\): the coefficients are signed numbers,
and an unwritten coefficient is one. The distributive law gives
\[
 3(x+2)+4x=3x+6+4x=7x+6.
\]
Combining visible numbers while ignoring their scope instead produces
\(7x+2\), which already disagrees at \(x=0\), giving two rather
than six. A procedure that works on the first example is therefore not
established as a correct general procedure. This is a worked mathematical
contrast from the source's component account, not a measured instructional
gain \cite[galley p.~701]{KoedingerRoll2011TransferTypeset}.
\end{hbexample}

\begin{hbexample}{def:p19-cta-function-induction}{A strategy and the missing recognition step}
The source discusses induction from input/output pairs
\((2,1),(3,3),(4,6),(5,10)\). Dividing each output \(y\) by
its corresponding input \(x\) gives \(1/2,1,3/2,2\).
Recognizing increments of one half suggests
\(y/x=(x-1)/2\), hence the candidate \(y=x(x-1)/2\).
The division strategy and recognition of the resulting pattern are
different contributions.

The four observations do not logically force that function. For any
real parameter \(\lambda\), the function
\[
 g_\lambda(x)=\frac{x(x-1)}2+
       \lambda(x-2)(x-3)(x-4)(x-5)
\]
has the same four values, because one factor in the added term is zero
at each observed input. At \(x=6\), the added term is
\(24\lambda\). A task can have an intended rule or favor a simpler
candidate without turning finite interpolation into unique deduction.
The learner evidence cited by the galley concerns recognising a useful
candidate pattern; the nonuniqueness check here delimits the inference
\cite[galley p.~702]{KoedingerRoll2011TransferTypeset}.
\end{hbexample}

\section{How cognitive task analysis obtains evidence}
\label{sec:p19-cta-methods}

The source describes procedural production models with application
conditions and actions, subgoals or retrieval requests. Declarative
retrieval is also context-sensitive. These descriptions motivate
isolating a proposed component without discarding the cues that make
its use appropriate. A computational model's description of a mental
function is not the mental function itself; successful prediction does
not establish uniqueness of the modeled architecture
\cite[galley pp.~703--704,713]{KoedingerRoll2011TransferTypeset}.

\paragraph{Think-aloud analysis.}
Record what novices and experts do while solving specified tasks,
including intermediate representations, attempted operations, revisions
and impasses. Verbalized reasons and observable operations are different
evidence. A fluent expert may not articulate every acquired processing
detail. The source proposes also studying experts working outside their
familiar content domain, to avoid confusing domain knowledge with
domain-general procedures. Unreported thoughts cannot simply be filled
in from the analyst's preferred explanation.

\paragraph{Difficulty Factors Assessment (DFA).}
Construct matched problems varying specified factors, distribute them
across forms, and retain which items each learner encountered. Examples
include story versus word-equation versus symbolic presentation,
result versus start unknown, and whole-number versus decimal arithmetic.
A change in correctness across forms can locate a difficulty factor;
it does not independently identify the one internal component causing
that difficulty. Counterbalancing, scoring and actual item sampling
matter as much as the factor names.

\paragraph{Longitudinal tutor data.}
Step-level logs can record an initial response, an error, a hint, and
the eventually completed step. Assistance permits later steps to be
attempted even after an early error, and sequences reveal changes over
practice. Such logs are not automatically as balanced as deliberately
constructed DFA forms. Adaptive item selection, repeated exposure,
hint use and changing item difficulty can affect a learning curve.
First-attempt success and eventual supported success answer different
questions \cite[galley pp.~708--711]{KoedingerRoll2011TransferTypeset}.

The galley reports a seven-study review favoring instructional redesign
based on task analysis, with a mean effect size of 1.7. That is its
report of Clark and colleagues, not a new pooled analysis with one
participant denominator or a demonstrated universal gain. The source
does not supply the uncertainty and study-level heterogeneity needed
to turn that number into a prediction for a new curriculum.

\section{Matched representations and actual solution procedures}
\label{sec:p19-cta-representations}

Koedinger and Nathan distinguish three external forms: a contextual
story, a situation-free verbal description of operations, and a symbolic
equation. The comparison can separate a benefit of familiar linguistic
form from a benefit of the particular situation. Symbolic sentences
must themselves be understood: equality, the order of operations,
grouping, variables and argument order do not become transparent merely
because the sentence is concise
\cite[pp.~131--139]{KoedingerNathan2004StoryCTA}.

\begin{hbexample}{def:p19-cta-unwind}{A solved problem without an obligatory equation stage}
A source problem concerns a worker earning a total of 81.90 monetary
units after six hours of work plus 66 in tips. Let \(x\) be the
hourly wage in monetary units per hour. The symbolic relation is
\(x\times6+66=81.90\).

One correct procedure first removes the tips, obtaining
\(81.90-66=15.90\), then divides the remaining earnings by six,
giving \(15.90/6=2.65\). This \emph{unwind} procedure inverts the
last operation first. A symbolic procedure subtracts 66 from both sides
of the equality, obtaining \(x\times6=15.90\), and divides both
sides by the nonzero number six, obtaining \(x=2.65\).
Substitution checks the answer: \(2.65\times6+66=81.90\).

A guess-and-test procedure can also test candidate wages and compare
their resulting totals with 81.90; stopping at a value with the wrong
total does not solve the problem. These procedures share quantitative
constraints but need not have the same intermediate written form.
The existence of the unwind route refutes the premise that solving this
story necessarily requires first writing the displayed equation.
It does not prove that unwind is the better instructional strategy for
every learner or problem \cite[pp.~135,145--148]{KoedingerNathan2004StoryCTA}.
\end{hbexample}

The 2004 studies crossed presentation, unknown position, number type,
and final arithmetic. A result-unknown problem requests the output
of an already specified sequence; a start-unknown problem requests an
input consistent with a specified output. For a described operation
sequence, directly applying the forward operations cannot alone recover
the missing input. This procedural requirement is not a theorem that
every person must find every start-unknown problem harder than every
result-unknown problem, independent of their knowledge and strategies.

\subsection{The two high-school assessments}

For the empirical records below, \(N\) is a participant count, \(F\)
the reported analysis-of-variance statistic with its displayed degrees
of freedom, \(p\) its reported significance probability, and \(d\)
the reported standardized Cohen contrast. SD denotes a reported
standard deviation. These quantities are not the algebraic variables
in the worked examples; in particular, the effect-size \(d\) is
distinct from the monetary amount \(d\) below. A non-significant
test is not a demonstration of exact equivalence.

DFA1 included 76 students at one urban high school: 58 in three
mainstream Algebra 1 classes and 18 ninth-grade geometry students who
had taken Algebra 1 in eighth grade. DFA2 included 171 first-year
Algebra 1 students from 24 classrooms at three urban high schools.
These are two samples, not two conditions randomly assigned within
one training experiment.

Ninety-six problems derived from four cover stories crossed three
presentations, two unknown positions, two number types, and two final
arithmetic types. Sixteen forms contained eight problems each. Four
were stories, two word equations and two symbolic equations. Form
construction prevented repeated base numbers or answers within a form
and counterbalanced problem order. The four stories on a form used
one number type and the other four items the other number type.
DFA2 changed one decimal base equation from the first study. Students
had eighteen minutes in class, were asked to show work and mark a
final answer, and were instructed not to use calculators
\cite[pp.~139--142]{KoedingerNathan2004StoryCTA}.

The overall representation means were:
\begin{table}[htbp]
\centering
\begin{tabular}{lrrr}
\hline
Sample &Story&Word equation&Symbolic equation\\ \hline
DFA1, 76 students&66\%&62\%&43\%\\
DFA2, 171 students&70\%&61\%&42\%\\ \hline
\end{tabular}
\caption{Reported correctness by representation in the two 2004
assessments; these are not treatment gains.}
\label{tab:p19-cta-dfa}
\end{table}

Both item and participant analyses supported the main verbal/symbolic
contrast. The distinction between story and situation-free wording was
more restricted. In DFA2, whole-number story and word-equation
correctness was 72\% and 73\%; for decimals it was 68\% and 48\%.
The representation-by-number-type interaction was detected in DFA2
but not DFA1. Familiar situation and familiar language are therefore
not interchangeable explanations
\cite[pp.~142--145,Figs.~3--4]{KoedingerNathan2004StoryCTA}.

Written-work coding also distinguished strategies. On start-unknown
stories, unwind was used in 50\% of responses, compared with 22\%
on word equations and 13\% on symbolic equations. Symbolic equations
elicited more blank responses. Among uses of the identified strategies,
unwind and guess-and-test had higher observed success than symbol
manipulation. Those conditional success rates are not randomized
strategy effects: learners chose strategies on potentially different
problems. The authors explicitly identify a choice/no-choice design
as a stronger test of strategy efficacy
\cite[pp.~145--150,Tables~4--5]{KoedingerNathan2004StoryCTA}.

Their error analysis separates blank responses, arithmetic errors and
conceptual errors; conditional error proportions use different
denominators. A manipulation on both sides of an addition sign, rather
than both sides of an equality, is not a valid equation transformation.
Adding a number of hours to an amount of money violates the quantities'
roles. In the symbolic form, these errors can remain less conspicuous
than in a meaningful story. The authors' comprehension explanation is
supported by the pattern of errors; a blank response alone does not
uniquely diagnose a missing parse rule
\cite[pp.~149--156,Figs.~5--6]{KoedingerNathan2004StoryCTA}.

\section{When abstract representations provide an advantage}
\label{sec:p19-cta-complexity}

Koedinger, Alibali and Nathan subsequently tested a
representation-complexity tradeoff, not a blanket advantage of stories
\cite[pp.~366--397]{KoedingerAlibaliNathan2008TradeoffsCTA}.
Their theoretical comparison distinguishes familiarity and semantic
redundancy, which can support grounded representations, from concise
externalization and manipulation, which can support abstract forms.
These are proposed mechanisms; they were not each separately randomized.

\begin{hbexample}{def:p19-cta-double-reference}{Combine constraints before reversing operations}
For the lottery example, let \(X\) be the prize amount in monetary
units. Keeping 64 and sharing the rest among three recipients, each
receiving 20.50, gives \((X-64)/3=20.50\).
Multiplying by three gives \(X-64=61.50\), and adding 64 gives
\(X=125.50\). At each inverse step one input and the result of the
relevant relation are known.

For the discount example, let \(X\) instead be the original price.
A reduction of 15\% of that price gives
\(X-.15X=38.24\). The same unknown contributes both the original
amount and the reduction. The equation reduces to
\(.85X=38.24\), hence \(X=38.24/.85\).
Taking 15\% of 38.24 and adding it does not invert the discount,
because the reduction was computed from \(X\), not from 38.24.
The exact-real answer is approximately 44.988235; obtaining a unique
cent-denominated price would additionally require a rounding convention.

The one-relation-at-a-time unwind procedure cannot proceed on the
initial two-unknown relations in the discount network. An informal
solver can nevertheless \emph{untangle} it by recognizing that the
sale price is 85\% of the original price. Repeated reference does
not make every non-symbolic procedure impossible; it changes what
constraint combination is required
\cite[pp.~370--372,388,394 n.~3]{KoedingerAlibaliNathan2008TradeoffsCTA}.
\end{hbexample}

\subsection{The two college samples and their measures}

Experiment 1 included 153 state-university students, 43 in elementary
and 110 in intermediate algebra. They had twenty minutes in class
for a six-item assessment and were randomly assigned one of six forms.
Eighteen items crossed three external forms---implicit-operator story,
explicit-operator story and equation---with single versus double
reference, using three base stories at each complexity level.
Forms counterbalanced position and presented all six base stories.
Other properties differed between the single- and double-reference
sets as well; reference count was not varied while every other
complexity property was held fixed.

The paper reports a representation-by-complexity interaction in the
participant analysis, \(F(2,304)=27.46\), \(p<.0001\).
The item analysis gave \(p<.10\), not a conventional .05 result.
Implicit and explicit stories did not have a detected overall
difference. The item-level pattern was not a uniform failure on all
equations: two familiar single-reference equation forms were answered
almost as well as stories, whereas the grouped division form was
substantially harder \cite[pp.~372--378]{KoedingerAlibaliNathan2008TradeoffsCTA}.

Experiment 2 included 65 Carnegie Mellon undergraduates receiving
course credit, with mean self-reported mathematics SAT score 719
(range 510--800). It included result-unknown and start-unknown
single-reference items as well as the three double-reference cases.
Each participant received story and equation versions using different
number sets. The interaction occurred in participant and item analyses.
Table~\ref{tab:p19-cta-complex-items} preserves the two original
samples' reported item contrasts rather than pooling their different
backgrounds and test forms.

\begin{table}[htbp]
\centering
\small
\begin{tabular}{p{.43\textwidth}rrrr}
\hline
 &\multicolumn{2}{c}{Experiment 1}&\multicolumn{2}{c}{Experiment 2}\\
Double-reference task &Story&Equation&Story&Equation\\ \hline
38 pupils, six more girls than boys&54\%&71\%&86\%&92\%\\
Price after a 15\% discount&4\%&29\%&63\%&75\%\\
Two exchange offers with different rates/fee&5\%&41\%&82\%&94\%\\ \hline
\end{tabular}
\caption{Reported correctness on the double-reference cases,
original Tables 5 and 7. Experiment 1 story percentages average its two
wording versions. Experiment 2 used paired forms with different numbers.}
\label{tab:p19-cta-complex-items}
\end{table}

Coding distinguished algebra, forward arithmetic, unwind, untangle,
guess-and-test, and answer-only responses. Formal/informal strategy
success remained an observational comparison, including analyses
restricted to people who used both. The proposed direct test requiring
one strategy or the other was not performed in these experiments.
For Experiment 1, a second trained coder rescored fifteen participants:
reported agreement was 96\% on ninety strategy judgments, and 93\%
and 84\% on forty-six error judgments at the broad and finer grains.
For Experiment 2, agreement was 96\% on one hundred strategy judgments
from ten participants. An additional thirty incorrect solutions yielded
93\% and 85\% agreement at the two error-coding grains.
These are agreement percentages under the specified coding samples,
not chance-corrected reliability or evidence that every hidden process
was directly observed.
The paper did not collect solution times, so efficiency is a
theoretical explanation rather than an independently measured speed gain
\cite[pp.~374,377--388,390--394]{KoedingerAlibaliNathan2008TradeoffsCTA}.

The contrast is tied to which task demands defeat which procedures.
The authors expressly note that some single-reference equations were
harder than some double-reference equations. Two occurrences of a
letter are not a universal difficulty metric, and a more concise
external representation is not automatically better understood.
Their proposed developmental sequence is a model consistent with the
contrasts, not a longitudinal observation of every learner following
one inevitable sequence.

\section{Algebraic production and the hidden component hypothesis}
\label{sec:p19-cta-production}

The distinction between solving a story for a number and producing an
expression for all relevant values changes the knowledge demand.
Producing a correct numerical instance does not alone establish
symbolization. Nor does recognizing one operation ensure that the
learner can embed its expression inside another.

The related galley names the proposed recursive-expression-production
component \texttt{RecExprProd}; the 2010 paper analyzes the corresponding
ability to embed one expression inside another. The galley contrasts
that component with \texttt{EngToMathOps}, comprehension of the
story's English and translation of its quantitative relations.
These are cognitive-model labels, not empirical proof that exactly
two indivisible mental modules exist.

The source's schematic expression grammar admits an expression
consisting of a quantity, an operator and another quantity, and extends
this to a quantity combined with an expression, or an expression combined
with a quantity. In the galley's wider hierarchy, \texttt{Exp} names
the expression category, \texttt{Op} an operator category, and
\texttt{Num} and \texttt{Var} numbers and variable tokens.
An \texttt{Exp} can be a \texttt{Num}, a \texttt{Var}, a parenthesized
\texttt{Exp}, or an \texttt{Exp Op Exp}; \texttt{Op} can be
\texttt{+}, \texttt{-}, \texttt{*}, or \texttt{/}.
Numbers, variables and grouped or recursively combined expressions
thus belong to the same expression category. These schematics
describe recursive structure; they do not by themselves supply an
unambiguous parser, precedence convention, semantic interpretation,
or rule for applying a formula in a story
\cite[pp.~471--473,Fig.~1]{KoedingerMcLaughlin2010LanguageTransfer}
\cite[galley pp.~707--710,Fig.~40.1]{KoedingerRoll2011TransferTypeset}.

\begin{hbexample}{def:p19-cta-students-substitution}{Compose two recognized relationships}
Let \(b\) be the number of boys, \(f\) the nonnegative number by
which boys are fewer than 62 girls, and use whole-number counts with
\(0\leq f\leq62\). The total is \(62+b\), and the boys count
is \(62-f\). Replacing \(b\) by the whole latter expression gives
\[
 62+(62-f)=124-f.
\]
The source also accepts the ungrouped spelling \(62+62-f\), under
ordinary arithmetic conventions. The answer \(62-f\) gives only
the boys count and omits the girls; \(62+b\) alone leaves the
requested expression dependent on an unspecified additional count.
This example distinguishes recognizing each relation from producing
their composition. The source variables and its numerical relationships
are retained \cite[pp.~472--473,Tables~1--2]{KoedingerMcLaughlin2010LanguageTransfer}.
\end{hbexample}

\begin{hbexample}{def:p19-cta-grouped-substitution}{Preserve the substituted expression's scope}
Let \(d\) be an amount to be divided equally among four recipients.
Let \(m\) be the amount spent from an initial amount of 72.
Substituting \(72-m\) for \(d\) in \(d/4\) gives
\((72-m)/4\), not \(72-m/4\). For \(m=20\), the two values
are 13 and 67. The substitution is of an expression, not a deletion
of its internal grouping requirement.

Let \(h\) be an original travel duration in hours and \(t\) the
new duration in the same units. Substituting \(h-2\) for \(t\) in \(550/t\)
gives \(550/(h-2)\). It is defined only when \(h\ne2\);
the source travel interpretation requires \(h>2\) for positive
duration. Replacing it by \(550/h-2\) changes the computation,
not merely its typography. These are mathematical scope checks,
not observations of which mistake any new participant made.
\end{hbexample}

\section{Worked symbolization of the eight instructional task types}
\label{sec:p19-cta-worked}

The eight tasks in the 2010 paper's Table 1 request expressions, not
numerical answers for an assigned value. The reconstructions below
give each relationship and its admissible interpretation. The domain
conditions make the mathematics explicit; they are not claimed as
additional instructions administered in the experiment. In the source
spelling, \(*\) denotes multiplication. A rational constant such as
\(1/3\) is treated as a quantity when counting the task's principal
operators, rather than every printed slash being a separately trained
operation \cite[pp.~472--475]{KoedingerMcLaughlin2010LanguageTransfer}.

\paragraph{Deliveries.}
There are five loads, twelve boxes per load, and \(c\) CDs per box.
With \(c\) a nonnegative whole count, the number of boxes is
\(5*12=60\), so the total is \(60*c=5*12*c\).
The result combines both multiplicative relations; \(12*c\)
counts a single load rather than all five.

\paragraph{Net earnings.}
For \(h\) hours at five monetary units per hour, gross earnings
are \(5*h\). Subtracting the seven-unit taxi cost gives \(5*h-7\).
Here \(h\geq0\) is a duration, and the result can be negative
if the cost exceeds the earnings; nonnegative gross earnings do not
logically ensure nonnegative net earnings.

\paragraph{Pooling and sharing.}
One parent has saved 972 monetary units and the other \(b\geq0\).
The pooled amount is \(972+b\); equal division among five children
gives \((972+b)/5\). In \(972+b/5\), only the second parent's
amount is divided. At \(b=28\), equal sharing gives 200, not 977.6.

\paragraph{Spending and sharing.}
An amount 72 is reduced by a mother's present costing \(m\),
leaving \(72-m\). Sharing that remainder equally among four
sisters gives \((72-m)/4\). The savings interpretation permits
\(0\leq m\leq72\). Both the subtraction and the scope of the
division are needed, as in the preceding substitution example.

\paragraph{Pupils.}
The girls count is 62 and the boys count is \(62-f\), giving
\(62+(62-f)=124-f\). Reusing 62 in both relationships does not
introduce two different girls populations. A variable for the
shortfall is not the boys count itself.

\paragraph{Approaching a dock.}
Initially 800 yards away, the rower approaches at forty yards per
minute. For \(m\) minutes the travelled distance is \(40m\),
and before or at arrival the remaining distance is \(800-40m\),
with \(0\leq m\leq20\). Beyond that interval, the linear
expression becomes negative; extending it unchanged as an ordinary
nonnegative distance would require a different motion description.
The time variable is \(m\), as in the original 2010 task.

\paragraph{Faster travel.}
A 550-mile journey normally taking \(h\) hours is completed two
hours sooner. The new duration is \(h-2\), not \(h\) and not
\(550/h\). For \(h>2\), average speed is distance divided by
duration: \(550/(h-2)\) miles per hour. At \(h=7\), this is
110; the ungrouped expression \(550/h-2\) gives a different value
and also subtracts a bare time change from a rate.

\paragraph{A one-third reduction.}
For original price \(d\geq0\), the reduction is \((1/3)*d\),
so the remaining price is \(d-(1/3)*d=(2/3)*d\).
The task uses \(d\) in both relations. A correct simplified
one-multiplication answer is possible even though the unsimplified
task combines two relations. Counting printed operators therefore
cannot substitute for checking semantic correctness.

\section{The 2010 substitution intervention}
\label{sec:p19-cta-intervention}

\subsection{Population, assignment and item sequence}

The classroom study used ASSISTment, which supplies feedback, optional
hints and scaffolding questions. During the 2008--2009 school year,
318 middle-school students in an urban district outside Boston
provided initial records: 158 in one-step-story practice and 160 in
substitution practice. The analysis retained the 303 completing all
sixteen items, \(N=303\), comprising 154 and 149 respectively. The six-page original method
does not explicitly describe the randomization mechanism. A
comparison of the supplied conditions should not silently invent its
implementation, nor should similar pretest means prove exchangeability
\cite[pp.~473--474]{KoedingerMcLaughlin2010LanguageTransfer}.

The What Works Clearinghouse's May 2012 review separately classifies
this study as a randomized controlled trial meeting its standards
without reservations under review protocol and standards 2.1.
That is an attributed design qualification, not an assignment procedure
described in the six-page report. The review also states that its
rating concerns this study rather than the entire evidence base for
the intervention \cite{WWC2012KoedingerMcLaughlinCTA}.

Thirty-two possible items were generated from eight two-step stories,
their paired one-step versions, and substitution versions. Each learner
saw sixteen: four two-step, eight one-step and four substitution items.
Two versions avoided repeated answers or answer parts; each had a
reversed difficulty order. Five initial assessment items comprised
four one-step items and one two-step item. Seven interleaved
instruction/assessment items followed, then four fillers.

In the source sequence codes, \texttt{a,b} are the two forms
of one-step items, \texttt{c} a substitution item, and
\texttt{0--3} two-step items in the assigned difficulty order.
The one-step condition's sequence is the source's
\texttt{abab0; ab1a2b3}, while the substitution condition's is
\texttt{abab0; cc1c2c3}.
The semicolon separates initial and interleaved phases; the four
later fillers are not represented in this abbreviation.
All items could subsequently provide assistance, including assessment
items. This is not a clean sequence of unassisted pretest, isolated
instruction and delayed unassisted transfer test.

The pretest score weighted the first two-step item equally with
the average of four one-step items. The posttest score averaged the
remaining three two-step items. Scores used the first attempt;
an incorrect entry or hint request counted as an error.
That criterion is distinct from eventual correct completion after
help. The source supplies no age/sex distribution, language-history
distribution, delayed retention result or representative sampling frame.

\subsection{Measured contrast and errors}

The reported raw pretests were .56 for substitution and .57 for
one-step practice. Raw posttests were .39 (SD .35) and .33 (SD .33).
With pretest and item order included, adjusted posttests were .39
and .32, and the reported standardized contrast was \(d=.29\).
The condition analysis gave \(F(1,299)=4.45\), \(p<.05\).
Order had a strong effect: an easy-to-hard item sequence necessarily
leaves harder tasks for the later outcome, unlike its reverse.
These raw, adjusted and item-order comparisons are not interchangeable
effect estimates \cite[p.~474]{KoedingerMcLaughlin2010LanguageTransfer}.
The original report prints no confidence interval for the selected
condition contrast. Its standard deviations describe score dispersion,
not an interval around the adjusted difference.

The four error classes were one-operator, two-operator,
missing-parentheses, and hint/other. One-operator errors occurred
in 34\% versus 30\% of responses in the one-step and substitution
conditions. On the three parentheses-requiring tasks, the
missing-parentheses proportions were 25\% and 13\%.
The paper also examines selected item groups and alternative correct
forms. Its simplified correct answers include \(60c\), \(124-f\)
and \((2/3)*d\); correctly using redundant parentheses is another
form variation. Neither a response's visual length nor its operation
count alone decides correctness.

Both groups were 72\% correct on the one-step items, though the
substitution group's items occurred in the later filler phase and
the other group's in instruction. This constrains a simple
interpretation that substitution merely teaches indiscriminate
two-operator responding, but does not isolate a unique grammar-learning
mechanism. Scaffolded exposure to both response types, expression
fluency, grouping knowledge, and response-frequency effects remain
relevant explanatory possibilities
\cite[pp.~474--476]{KoedingerMcLaughlin2010LanguageTransfer}.

The transfer is from specified substitution exercises to different
story-symbolization tasks. It is not transfer to general proof,
unmeasured everyday reasoning or an unrelated subject. The authors
propose that producing, inspecting and receiving feedback on a
uniform symbolic answer can support abstraction even when source
and target prompts look dissimilar. The experiment did not separately
randomize output uniformity, perceptual recoding and each kind of
feedback. A successful instructional package does not by itself
establish which proposed mediator caused its difference.

\section{Overlapping components and alternative models}
\label{sec:p19-cta-overlap}

The galley contrasts an overlap model, in which task A requires
\texttt{K1} and task B requires \texttt{K1} plus \texttt{K2},
with a nonoverlap model assigning separate \texttt{Ka} and
\texttt{Kb}. In its example, recursive expression production is
the proposed shared component and story-to-operation comprehension
the additional one. Initial success on A versus B does not settle
which model is correct. Sequential or intervention evidence can
test whether successful practice on A helps B, provided order,
item difficulty, exposure and assignment are accounted for.

The galley's 38\%-versus-27\% comparison and accompanying note
describe an alternative analysis of the 2010 data. They must not
replace the original paper's aggregate raw .39/.33 or adjusted
.39/.32 outcomes. Even a transfer contrast supporting shared
knowledge does not identify the unique size or indivisibility of
that knowledge component
\cite[galley pp.~710,713 n.~9]{KoedingerRoll2011TransferTypeset}.

\section{Instructional redesign without universal prescriptions}
\label{sec:p19-cta-redesign}

The source advocates testing the actual target component and the
grain size of the selected analogies. Changing a plumber's rate
and fixed fee to an electrician's changes numbers and cover story
while preserving a positive-rate affine form. A negative rate,
division by an expression, repeated reference, and a required
grouping each add different demands. Equal success on one pair
does not demonstrate transfer across every broader class.

Its discussion of worked examples versus practice retains a
content-by-method hypothesis: example study might particularly
support complex schema induction, while earlier retrieval practice
might support simpler factual retention. The galley explicitly
identifies this interaction as untested to the authors' knowledge;
it is not a universal schedule validated by the algebra experiment.

Language-mediated and non-language-mediated learning are also
different proposed routes. The source uses language broadly for
external symbolic systems, while distinguishing verbal mediation
of learning from learning a symbolic pattern through examples or
feedback. Absence of an explicit spoken grammar lesson does not
show absence of covert language, awareness, or self-explanation.
Its open questions concern the routes' mechanisms, continued use,
and independence or interaction. No single inevitable order from
verbal understanding to practice, or the reverse, is established
\cite[galley pp.~705,707--713,nn.~6,10]{KoedingerRoll2011TransferTypeset}.

\begin{hbmethod}{def:p19-cta-redesign}{An actionable task-analysis cycle}
Specify the intended response and its correct interpretation, including
domains, units and application conditions. Solve representative
instances by more than one possible route before assuming that a
particular formal intermediate step is necessary.

Contrast item factors that can distinguish missing comprehension,
production, composition, retrieval or strategy knowledge. Preserve
the participant and item sampling, counterbalancing, response
opportunities, scoring denominators and assistance. Use observed
errors to formulate competing models rather than assigning a unique
component solely from a final incorrect answer.

Design an instructional contrast targeting the proposed missing
component, with a meaningful comparison condition. Specify in advance
which new tasks would constitute transfer and which would merely
repeat the trained form. Account for feedback and changing difficulty;
record first attempt separately from eventual supported completion.

Assess the complete intended outcome. A correct expression does
not show how it was produced; an instructional gain does not show
durable retention; and a broad component label does not prove
generality. If the redesigned activity improves one target but fails
when sign, grouping or domain changes, revise the component's
application conditions rather than silently classifying those cases
as outside algebra.
\end{hbmethod}

Chapter~\ref{sec:p19-self-explanation-and-learning} supplies the
distinct self-explanation protocols and their elicitation/coding
boundaries; its biology and diffusion samples are not participants
in these algebra studies. Chapter~\ref{sec:p19-analogy-and-problem-transfer}
supplies retrieval, noticing and adaptation comparisons. The
practice/transfer/retention distinctions in
\cref{def:p19-r153-transfer} delimit any broader training claim.
These interfaces clarify the task, but do not substitute those
sources' evidence for the present cognitive-task-analysis account.
